% ============================================================
%  ANÁLISIS NUMÉRICO O COMPUTACIONAL
%  Fancho → Actividad 8  (Entregar)
% ============================================================

\section{Análisis Numérico o Computacional:}

% ──────────────────────────────────────────────
%  ACTIVIDAD 8  →  Fancho  (Entregar)
% ──────────────────────────────────────────────

\actividad{8} \textbf{(Entregar).} [Mix Numérico -- No Rutinario] Considere el problema
con valor inicial:

\[
    \frac{dy}{dt} + \left(\sqrt{1+\sin^2(t)}\right)y = t, \qquad y(0)=2
\]

\begin{enumerate}

    % ── Punto 1 ─────────────────────────────
    \item Utilice la definición de integral o el Teorema Fundamental del Cálculo para
    deducir que el factor integrante para esta ecuación diferencial se puede escribir
    como:
    \[
        \mu(t) = e^{\int_0^t \sqrt{1+\sin^2(s)}\,ds};
    \]
    y que la solución del problema con valor inicial es:
    \[
        y(t) = \frac{1}{\mu(t)}\int_0^t s\,\mu(s)\,ds + \frac{2}{\mu(t)}.
    \]

    \begin{respuesta}[title=Punto 1 — Fancho]
        % Demostración analítica del factor integrante y la fórmula de solución
        \vspace{10cm}
    \end{respuesta}

    % ── Punto 2 ─────────────────────────────
    \item Obtenga una aproximación de la solución en $t=1$. Utilice un método de
    integración numérica (como la regla de Simpson) en un ciclo anidado para estimar
    valores de $\mu(t)$, y con ello, estimar el valor de $\displaystyle\int_0^1 s\,\mu(s)\,ds$.

    \smallskip
    \noindent\textit{[Sugerencia]}: primero utilice el método de integración numérica para
    estimar $\mu(t)$ en $t=0.1, 0.2, 0.3,\ldots,1$. Luego, utilice estos valores y
    aplique de nuevo el método de integración numérica para aproximar
    $\displaystyle\int_0^1 s\,\mu(s)\,ds$.

    \begin{respuesta}[title=Punto 2 — Fancho]
        % Desarrollo numérico: ciclo anidado con regla de Simpson
        % Incluir tabla de valores de mu(t) para t = 0.1, 0.2, ..., 1
        % y el valor aproximado de la integral y de y(1)

        \textbf{Tabla de valores de $\mu(t)$:}
        \vspace{0.3cm}

        \begin{center}
        \begin{tabular}{cccc}
            \toprule
            $t$ & $\sqrt{1+\sin^2(t)}$ & $\int_0^t\sqrt{1+\sin^2(s)}\,ds$ (aprox.) & $\mu(t)$ (aprox.) \\
            \midrule
            0.1 & & & \\
            0.2 & & & \\
            0.3 & & & \\
            0.4 & & & \\
            0.5 & & & \\
            0.6 & & & \\
            0.7 & & & \\
            0.8 & & & \\
            0.9 & & & \\
            1.0 & & & \\
            \bottomrule
        \end{tabular}
        \end{center}

        \vspace{0.5cm}
        \textbf{Estimación de $\displaystyle\int_0^1 s\,\mu(s)\,ds$:}
        \vspace{3cm}

        \textbf{Valor aproximado de $y(1)$:}
        \vspace{2cm}
    \end{respuesta}

    % ── Punto 3 ─────────────────────────────
    \item Utilice el método de Euler para aproximar la solución en $t=1$, con tamaños de
    paso $h=0.1$ y $h=0.01$.

    \begin{respuesta}[title=Punto 3 — Fancho]
        % Método de Euler con h = 0.1
        \textbf{Euler con $h = 0.1$:}
        \vspace{0.3cm}

        \begin{center}
        \begin{tabular}{ccc}
            \toprule
            $n$ & $t_n$ & $y_n$ \\
            \midrule
            0 & 0.0 & 2.0 \\
            1 & 0.1 & \\
            2 & 0.2 & \\
            3 & 0.3 & \\
            4 & 0.4 & \\
            5 & 0.5 & \\
            6 & 0.6 & \\
            7 & 0.7 & \\
            8 & 0.8 & \\
            9 & 0.9 & \\
            10 & 1.0 & \\
            \bottomrule
        \end{tabular}
        \end{center}

        \vspace{0.5cm}
        \textbf{Euler con $h = 0.01$:} \quad (mostrar solo resultado en $t=1$)
        \vspace{2cm}
    \end{respuesta}

    % ── Punto 4 ─────────────────────────────
    \item Como ingeniero: proponga una metodología para comparar la precisión y la
    eficiencia de los esquemas numéricos en 2 y 3.

    \begin{respuesta}[title=Punto 4 — Fancho]
        % Análisis comparativo: precisión (error relativo) y eficiencia (costo computacional)
        % Incluir: tabla comparativa de resultados, criterios de comparación propuestos,
        % conclusión sobre cuál método es preferible y por qué.

        \textbf{Tabla comparativa de resultados en $t=1$:}
        \vspace{0.3cm}

        \begin{center}
        \begin{tabular}{lcc}
            \toprule
            Método & $y(1)$ aprox. & Observaciones \\
            \midrule
            Simpson (Punto 2)   & & \\
            Euler $h=0.1$       & & \\
            Euler $h=0.01$      & & \\
            \bottomrule
        \end{tabular}
        \end{center}

        \vspace{0.5cm}
        \textbf{Análisis y conclusión:}
        \vspace{6cm}
    \end{respuesta}

\end{enumerate}

\newpage
